\documentclass{standalone}
\usepackage{circuitikz}
\usepackage[scaled]{helvet}
\usepackage[T1]{fontenc}
\renewcommand\familydefault{\sfdefault}


\begin{document}

\begin{circuitikz}[anchor=center]
    
\draw (0,0) node[dipchip,
    num pins=16,
    external pins width=0.6,
    scale = 2,
    external pad fraction=4,
    align=center ](C){\footnotesize BMI088\\ \footnotesize Break \\\footnotesize Out};
    
\draw[-latex] (C.pin 1) -- ++(-5,0) node[inner sep=0](SDO1){} -- ++(-1,0) node[anchor=east] {MISO};
\node [left,anchor=east] at (C.bpin 1) {SDO1};

\draw[latex-] (C.pin 2) -- ++(-1,0) node[anchor=east] {CSB\_ACCEL};
\node [left,anchor=east] at (C.bpin 2) {CSB1};

\draw[-latex] (C.pin 3) -- ++(-1,0) node[anchor=east] {INT4\_GYRO};
\node [left,anchor=east] at (C.bpin 3) {INT4};

\draw[-latex] (C.pin 4) -- ++(-1,0) node[anchor=east] {INT3\_GYRO};
\node [left,anchor=east] at (C.bpin 4) {INT3};

\draw[latex-] (C.pin 5) -- ++(-4,0) -- ++(0,1) -- ++(-0.25,0) -- ++(0.5,0) node[pos=0.5,anchor=south]{VDDIO};
\node [left,anchor=east] at (C.bpin 5) {VDDIO};

\draw(C.pin 6) -| (SDO1) to[short,-*] (SDO1);
\node [left,anchor=east] at (C.bpin 6) {SDO2};

\draw[-latex] (C.pin 7) -- ++(-1,0) node[anchor=east] {MOSI};
\node [left,anchor=east] at (C.bpin 7) {SDA/SDI};

\draw[-latex] (C.pin 8) -- ++(-1,0) node[anchor=east] {SCK};
\node [left,anchor=east] at (C.bpin 8) {SCL/SCK};

\draw[latex-] (C.pin 9) -- ++(4,0) node [inner sep=0](GND){} -- ++(0,-0.5) node[ground] {};
\node [right,anchor=west] at (C.bpin 9) {PS};

\draw (C.pin 10) -| (GND) to[short,-*] (GND);
\node [right,anchor=west] at (C.bpin 10) {GNDIO};

\draw[latex-] (C.pin 11) -- ++(1,0) node[anchor=west] {CSB\_GYRO};
\node [right,anchor=west] at (C.bpin 11) {CSB2};

\draw (C.pin 12) -| (C.pin 10 -| GND) to[short,-*] (C.pin 10 -| GND);
\node [right,anchor=west] at (C.bpin 12) {GNDA};

\draw[latex-] (C.pin 13) -- ++(3,0) -- ++(0,0.5) -- ++(-0.25,0) -- ++(0.5,0) node[pos=0.5,anchor=south]{VDD};
\node [right,anchor=west] at (C.bpin 13) {VDD};

\draw (C.pin 14) -| (C.pin 12 -| GND) to[short,-*] (C.pin 12 -| GND);
\node [right,anchor=west] at (C.bpin 14) {``NC''};

\draw[latex-latex] (C.pin 15) -- ++(1,0) node[anchor=west] {INT2\_ACCEL};
\node [right,anchor=west] at (C.bpin 15) {INT2};

\draw[latex-latex] (C.pin 16) -- ++(1,0) node[anchor=west] {INT1\_ACCEL};
\node [right,anchor=west] at (C.bpin 16) {INT1};

\end{circuitikz}

\end{document}
